\renewcommand{\thechapter}{1}
\chapter{Introduction}
\label{chapter:10-Introduction}

\section{Motivation}
Graphs are expressive data structures capable of encoding evidence as vertices and contextual relationships between them as edges. 
As powerful abstractions of the real-world, graphs lend themselves well to learning and reasoning tasks, which are encompassed in an emerging area of work we term Graph-Based Artificial Intelligence (\gai).
We taxonomize \gai within the existing AI landscape, covering the well-known areas of \textit{Graph Prediction}~\cite{Smith2021, Zhou2020a} and \textit{Graph Mining}~\cite{Chen2022a}, and extending those to include \ac{GPS}, \ac{GR}, \ac{GA}, and \ac{GT}, which are detailed in Figure~\ref{figure:gai-taxonomy}.
Across these areas, \gai includes tasks like structural analysis and inference over graphs, as well as applications of graphs to AI problems like model transparency, explainability, and interpretability.

\nrscomment{dropped this citation "specific problems like signal processing~\cite{Xia2021}"}
\nrscomment{changed Smith2021 and Zhou2020a from being GNN to broader GP, so as to not list GP redundantly in the list of things you extend with}
    
    \begin{figure}[h!]
        \begin{center}
            \includesvg[width=\columnwidth]{SVGs/gai_taxonomy}
            \caption{The \gai Taxonomy with example tasks}
            \label{figure:gai-taxonomy}
        \end{center}
    \end{figure}



Although graphs are expressive and broadly applicability to a variety AI problems, they impose a significant computational burden on \gai systems.
\gai algorithms that operate on graphs frequently have quadratic or worse complexity, which limits their usability on real-world graphs containing billions or trillions of vertices. 
In many non-graph settings, modern general-purpose processors use techniques like caching and prefetching to speed up computation over large inputs.
The challenge with graphs is that most \gai algorithms exploit the graph's structure in processing, which requires access to the data in topological order.
That requirement conflicts with typical graph storage, which is chronological in the order that vertices are created, meaning topologically close vertices are not necessarily proximal in memory. 
Without caching benefits, \gai algorithms rely on data retrieved from main memory, which causes memory latency and bandwidth bottlenecks that limit scalability and prevent \gai from being deployed on real-world problems that might otherwise benefit from it. 

\section{Problem and Background}
\label{section:11-introduction-to-problem}

Common approaches to speeding up \gai processing are narrow, and typically involve hardware, like specialized accelerators or distributed processing, representation, like custom data structures or abstraction models, or algorithms, like parallel or approximate methods.
These approaches are interdependent, meaning decisions about optimizing hardware can limit choices about parallelism or programming models, which makes it difficult to determine the optimal configuration for a given collection of problems. 

This thesis aims to determine \textit{if}, \textit{how}, and \textit{when} to speed up problem-solving on \gai given a set of operating \textit{constraints}.
Specifically, we consider how to effectively combine representations, hardware accelerators, and algorithms to speed up \gai, and under what conditions a specific combination is appropriate for a use case given some constraints on a solution's efficiency, effectiveness, and cost.

Prior approaches to identifying if, how, and when to accelerate a \gai system rely on generic competitive benchmarks or invasive workload analysis, which are both problematic. 

\textit{Benchmarking} is a process by which systems or components are compared according to some criteria, which is used to determine their relative performance~\cite{Kistowski2015}.
Most benchmarks for \gai take a `black-box' approach to competing systems based on \textit{efficiency} or \textit{effectiveness} metrics but rarely both, and never their interactions.
Benchmarks typically focus on evaluating isolated \textit{kernel} level tasks rather than more complicated real-world problems and typically use a small selection of datasets that do not cover the breadth of expected input sizes and shapes.
Existing benchmarks also lack support for dynamic graph processing approaches, especially comparing dynamic and static approaches under the same evaluation framework.

Alternatively, workload analyses are `white-box' approaches that analyze the internal state of a system while completing a workload to identify bottlenecks. 
Although workload analyses yield detailed information about how to accelerate the system, they also require access to the system under test for profiling and access to the code for instrumentation. 
The invasive configuration requirements mean workload analyses cannot be scaled for general-purpose comparisons, making them unsuitable for broad application to \gai.

In addition to the choice in type of analysis, several other challenges underlay the broader problem of deciding if, how, and when to improve \gai under some constraints, which we outline in Figure~\ref{figure:gai-challenges}.
Firstly, the user must decide the degree to which effectiveness should be traded for efficiency or cost.
The need to make this trade-off stems from the real-word constraints imposed on \gai systems, such as a requirement to generate a recurring solution at a pace more frequent than the runtime of the algorithm allows.
Specialized \ac{ASIC} hardware accelerators for graph processing can be optimized to speed up a specific algorithm, representation, abstraction, and programming model, but a given \gai system may need to support many graph tasks, so a specialized accelerator may not garner a broad enough improvement to justify its cost investment. 
These real-world challenges and trade-offs are particularly pronounced for \gai because of the inherent nature of graph processing that underlies \gai problems.
Graphs of different shapes and sizes influence algorithmic complexity in unpredictable ways, including by exacerbating the irregular memory accesses and computation patterns that undermine spatial and temporal locality. 
There is no standard accepted graph representation method, abstraction model, or set of primitive operations, which further prevents a unity of effort towards accelerating graph processing broadly. 


\begin{figure}[h!]
    \begin{center}
        \includesvg[width=\columnwidth]{SVGs/gai_challenges}
        \caption{\gai challenges are a hierarchy of dependencies foundationally limited by graph processing.}
        \label{figure:gai-challenges}
    \end{center}
\end{figure}




\section{Methodology}
\label{section:graph-based-ai-problem-methodology}
% \begin{landscape}
    
    \begin{figure}[h!]
        \begin{center}
            \includesvg[width=\columnwidth]{Figures/SVGs/benchmark_design_process}
            \caption{\gai Benchmark Design Framework}
            \label{figure:framework}
        \end{center}
    \end{figure}

% \end{landscape}

We address these \gai issues by contributing a methodology for cunstructing a problem-centric, accelerator-agnostic `grey-box' \gai benchmark that can evaluate both the \textit{efficiency} and the \textit{effectiveness} of a system under test.
We employ a formal method (Figure~\ref{figure:framework}) that emphasizes the transparency, verifiability, and reproducibility of benchmark problems and allows for extension to include emerging \gai problems. 
Our method comprises six steps that generate the components of the benchmark specification and the observation apparatus.

The process begins by identifying a real-world use case to ground the problem in real-world datasets and current approaches. 
A more detailed analysis of the problem yields information about the size and shape of the data and formally specifies the everyday tasks associated with that problem in natural language. 
Formalizing the study's phenomena includes identifying what specific interactions of dataset and task are likely to yield interesting results, providing the first significant contribution to the problem's `grey box' approach. 
Considering the phenomena, exploratory analysis of the domain builds the base of knowledge to make initial hypotheses about behavior. 
The hypotheses drive the data collection to study the phenomena of interest, and deviation from those hypotheses indicates interesting behavior worthy of further study. 
The construction of the apparatus involves collecting and generating the datasets that will be used to trigger the conditions the hypotheses predict will stimulate the phenomena of interest. 
It also involves developing any required \ac{API} extensions to facilitate the problem and evaluate the solutions. 
The methodology for designing benchmark problems yields a collection of datasets, tasks, and metrics that will combine to provide a `grey box' view of the \ac{SUT} stimulating phenomena of interest according to the defined hypotheses. 



\section{Our \gai Benchmark Framework}
% a tool for converting \ac{OSM} data to directed, weighted, attributed graphs with ground-truth communities, a benchmark for the \gai problem of population segmentation and a collection of experiments to demonstrate how the \gai benchmarking framework supports engineers and procurement officers in identifying if, when and how to accelerate systems for \gai.
\nrscomment{this is currently jut a collection of statements complimenting and describing the benchmark. need to sort out what belongs in method vs here, and how to reorganize this section to present a coherent picture of how we apply the method to construct something useful.}

Following our proposed methodology, we construct a \gai benchmark that balances the generality of competitive benchmarks and the detail of workload analyses without requiring invasive system access. 
We leverage constraint theory with an empirical approach to algorithmic analysis to develop inputs that specifically trigger bottleneck behavior. 
We then use our observations of those behaviors to develop a `grey-box' understanding of the system's internal state. 
Our approach is \textit{problem-centric}, meaning that the benchmark tasks are real-world problem statements rather than simplified kernels. 
For example, the \textit{kernel} task `\textit{run \ac{BFS}}' could be presented in problem-centric form as `\textit{Find the neighbors of vertex X}.'
The problem-centric view better captures the deployment of a system in the real world and encourages innovative approaches by not constraining the benchmark to specific approaches. 

The framework adopts an Input/Output that maximizes the degree of freedom of a benchmark user. 
Users can vary their hardware, software stack, and implementation without changing how their code interfaces with the benchmark. 
This flexibility allows for easy experimentation with different combinations of acceleration techniques to determine what works best for their use case. 
Finally, our evaluation approach explores the interaction between efficiency and effectiveness, recognizing that a fast partial solution can be more beneficial than a slow complete one.
Our more nuanced view of effectiveness makes it possible to evaluate dynamic and static approaches using the same framework.
Our repeatable process delivers an extensible framework that can be continually improved to evaluate new \gai problems as they emerge.

    
    \begin{figure}[h!]
        \begin{center}
            \includesvg[width=\columnwidth]{SVGs/gai_benchmark_high_level.svg}
            \caption{The high-level view of \gai benchmarking framework}
            \label{figure:gai-framework-overview}
        \end{center}
    \end{figure}



% The \gai Benchmark Framework shown in figure~\ref{figure:gai-framework-overview} is an Input/Output benchmark that evaluates the performance of a system under test by providing problems, recording the time to solution, and verifying the correctness of an answer. 
% The ideal system under test achieves a perfect effectiveness score in minimal time and cost, but all three dimensions will vary in the systems assessed. 

% The benchmark is designed so that the system under test is completely independent of the benchmark itself. 
% Using the Input/Output paradigm only requires the \ac{SUT} to interact with an \ac{API} (the observation apparatus). It so places no constraints on how a \ac{SUT} configures its hardware, software stack, or specific \gai implementation to solve the problem. 

% The benchmark definition of the \gai framework is comprised of independent benchmark problem modules. 
% Each module contains a collection of datasets designed to stress the efficiency and effectiveness of \ac{SUT} and uses \textit{problem-centric} task definitions to ensure problems are grounded in the real world and to encourage innovative solutions by removing the constraints on kernel-level tasks most existing benchmarks enforce. 

% The independence of the benchmark specification from the \ac{SUT} means that an analyst, engineer, or researcher using the benchmark can quickly modify their \ac{SUT} and validate performance changes on a known benchmark. 
% The careful design of the component benchmark problems to achieve the `grey-box' analysis of a \gai \ac{SUT} requires a formal methodology for transparency, reproducibility, and verifiability. 


\section{Impact}
The \gai Benchmark Framework is a flexible tool for experimentation and exploratory analysis for analysts, engineers, and researchers.
Analysts who routinely work with graph data and apply \gai to solve problems can use the \gai benchmark to understand better the sizes and shapes of inputs that will impact their work performance and use prior results to build estimates for \gai execution time. 
Engineers can use the \gai benchmark to evaluate which systems they should procure and to make architectural decisions about how their system implements \gai to maximize the efficiency and effectiveness across the breadth of tasks the system must support.
Researchers building \gai can use the \gai Benchmark Framework to experimentally vary the hardware and implementation details to find the bottlenecks in \gai as a function of the input datasets and real-world problems it solves. 
The ability to satisfy the needs of these three distinct user groups is a valuable feature of the \gai benchmark framework.


\section{Thesis Outline}
\label{section:15-Thesis-Outline}
The thesis continues with a review of the \gai taxonomy and key graph and computer architecture concepts in Chapter~\ref{chapter:20-Preliminaries} to identify what makes optimizing \gai approaches hard and why measuring relative performance is difficult. 
Then, Chapter~\ref{chapter:30-Literature-Review} is a detailed review of current \gai benchmarking and workload analysis efforts to characterize the current gaps in benchmark coverage of \gai tasks, datasets, and metrics and identify what causes performance to degrade in \gai systems. 
Chapter~\ref{chapter:40-Graph-AI-Benchmarking} introduces the \gai benchmark frameworks and outlines our methodology for developing \gai problems to constitute the benchmark framework. 
We use the population segmentation problem from the Intelligence and Security domain to demonstrate how our development methodology is applied, describe the implementation of the \gai benchmark prototype, and outline the findings of our initial experiments in chapter~\ref{chapter:50-Community-Detection-Case-Study}.
Finally, chapter~\ref{chapter:60-Conclusion} concludes after identifying specific directions for future work.